% 02-business-model.tex - Business model clarification and overview

\section{Business Model: Marketplace + CAC Optimization}

\subsection{Revenue Stream Segmentation}
\begin{itemize}
  \item \textbf{S1 --- Subscriptions (B2C marketplace, unit: individuals):} Anonymous marketplace for recurring memory-intensive computation services. \textbf{This is the only segment included in the Year 1--3 financial projections.}
  \item \textbf{S2 --- Licensing (B2B, unit: institutions):} Custodial TGPO licensed to banks, payment institutions and fintechs, where the \emph{institution} imposes the delay. Annual licence, seats and integration fees only --- \textbf{zero marketplace revenue and no \numint{\markPlatformFeePercent}\% take rate by construction}, because an institution that wants the delay to be real never buys solving capacity to shortcut it. Sized in the Market Opportunity section; excluded from projections.
  \item \textbf{S3 --- Capacity sales (B2B marketplace, unit: companies):} Solving capacity sold to companies protecting \emph{digital} assets (treasury keys, signing credentials, root secrets). Same take rate as S1, far higher ARPU, far lower count. Sized in the Market Opportunity section; excluded from projections.
  \item \textbf{Merchandise (Community Engagement):} Branded items sold via a separate 'Superfan Store' to engage our most passionate users. This is treated as a bonus revenue stream and is not included in the primary financial projections.
\end{itemize}

\noindent\textit{The three segments are counted in three different units and are never summed into a single customer number. Only annual revenue in dollars is comparable across them.}

\subsection{Subscription Service: Anonymous Computation Marketplace}
The subscription service operates as a computation matchmaking \textbf{marketplace} connecting:
\begin{itemize}
  \item \textbf{Clients:} Users needing recurring memory-intensive computation without owning adequate hardware - anonymity critical for privacy
  \item \textbf{Providers:} PC owners monetizing idle computational capacity - can operate publicly to attract clients
\end{itemize}
Platform provides anonymous client matching, reputation system, and dispute arbitration.

Key marketplace dynamics:
\begin{itemize}
  \item Computation jobs are simple but, by design, memory-intensive, lengthy and client demand for them is recurrent
  \item Anonymity is critical for Clients, while Providers can promote services openly
  \item We align provider success with platform growth. An aggressive referral commission incentivizes providers to recruit clients, transforming them into a highly-motivated, capital-efficient sales force \cite{parker2016}.
  \item First-mover advantage critical to build reputable user base before competitors
  \item Tiers differentiate by computation duration: \numint{\computeHoursBasic}, \numint{\computeHoursMedium}, \numint{\computeHoursProfessional} and \numint{\computeHoursGolden} hours
\end{itemize}

\begin{table}[H]
\centering
\caption{Customer Segment Characteristics}
\begin{tabularx}{\linewidth}{l X X}
\toprule
Attribute & Subscription Users & Free Users \\\midrule
Technical Level & Low-Medium & High \\
Purchase Preference & Recurring & Self-hosted (no subscription) \\
Price Sensitivity & Medium & High \\
SAM Size (S1 only) & \numint{\samSegOneIndividuals}~individuals\cite{henley2025crypto,chainalysis2024,triple2023} & \textemdash \\
\bottomrule
\end{tabularx}
\end{table}

\noindent\textit{SAM here refers to the S1 (B2C marketplace) Serviceable Available Market only, counted in \textbf{individuals}. The S2 institution count and the S3 company count are different units and are reported separately in the Market Opportunity section; they are never added to this figure. Free (self-hosted) users are excluded from SAM and related financial projections for simplicity. Merchandise is treated solely as community engagement and is not included in financial projections.}
